% !TeX root = ../papers/nc-p3.tex
\section{三つの長方形の水平成分と両符号の垂直成分を通る無制限変形}
\label{sec:rectangle-unrestricted-components}
\begin{theorem}\label{thm:rectangle-unrestricted-components}
二つの隣接切断がちょうど階数6のアフィンテンソル軌跡を $Z$ とする。
76型と112型の水平主成分、および第三長方形族の全閉包は、$Z$ 内で
定理\ref{thm:full-fifteen-component}の同じ既約成分に属する。
両型の $t=1$ の垂直族は、これと異なる一つの既約成分に属し、その一般点は両側の反対位置を
それぞれ対称化できる。両型の $t=-1$ の垂直族も一つのさらに異なる既約成分に属し、その成分は
定理\ref{thm:klein-other-pairings}の $\lambda=-1$ の表示の全基底変換による閉包である。
三成分のアフィン次元はそれぞれ $76,75,65$、指定始点の一般同型類の次元は $15,14,4$ である。
各々の幾何学的一般点は標準Koszul左右AS正則である。両符号の水平・垂直交叉はそれぞれ異なる次元の成分に属する。
\end{theorem}

\begin{lemma}\label{lem:finite-flat-conormals}
命題\ref{prop:twenty-conormal}の横断的二十点を、長さ20の有限完全交叉 $X$ に置き換えられる。
二つの部分評価には全台点で階数3を要求する。余法線の独立条件は、四配置の斉次座標の積
$a_i b_j c_k d_l$ の $X$ 上への制限が、対応する可逆層の全切断の二十次元空間を張ることである。
この条件下で二つのSchur微分の合計階数が180なら $Z$ はその点で滑らかな76次元である。
\end{lemma}
\begin{proof}
第一切断だけをその滑らかな階数6軌跡上で動かす。第一辺の六双一次式が定める射影族は
この点の近くで有限であり、相対完全交叉の平坦性判定で階数20の有限平坦射になる。
局所基底環上のその代数を $D$ とする。特殊ファイバーは有限個のArtin局所環の積なので、
各射影座標には全台で消えない一次形式が選べる。評価像の補空間 $H$ も同時に選べる。
命題\ref{prop:twenty-conormal}の $X=QJ_a$, $Y=J_b^{\mathsf t}P$ の議論を $D$ 上で行う
（この式の $X$ は行列を表す）。用いる三次小行列は各台点で選べ、核の直線束は
基底を局所化して自明化できる。これは特殊ファイバー上で生成元を選びNakayamaの補題で持ち上げればよい。
第二切断の階数方程式を課す前から、百個のSchur成分に対する同命題の恒等式が $D$ に値を持つ。
$D$ の自由基底でその係数を読めば基底環上の二十本の恒等式となる。
仮定した全切断の生成性は、まさに基点でこの二十行が独立であることを表す。
非主行・列への射影が左右核上の同型であることも各点での可逆小行列から従う。
従って先の命題と同様に八十本の生成系が得られ、微分階数により滑らかである。
\end{proof}

\begin{proof}[水平の三族を通る成分]
76型で $(u,v,z)=(2,3,5)$、112型で $(u,v,z,w)=(2,3,5,4)$、
第三長方形で $(A,D,b)=(2,3,5)$ を取る。素数1009で二切断のSchur微分を四文字に分けて計算すると、
共通接空間の文字別次元は三例とも $(18,18,20,20)$、基底変換からの次元は $(13,16,16,16)$ である。
第一辺を第一射影の四次小行列式で求め、一つのアフィン超平面で全台を覆う。
その座標環を標準単項式の二十次元空間とし、線型形式を可換な乗法行列として扱う。
各隣接評価の三行結合に規格化一次式を加えた四行四列の行列が、この環上で可逆である。
これで $a,b,c,d$ を順に解き、残りの評価式と一周後の一致、$a$ から逆に解いた $d$ も検算する。
したがって第一辺はこの有限スキームのグラフで、六式の有限完全交叉である。
四座標の積の256列は環の単位に作用させると階数20を持つ。
有限平坦な点環と単元小行列を素数で局所化した基底に持ち上げれば、標数零で補題が適用できる。
特にどの点も全 $Z$ の滑らかな点である。

四元群の固定軌跡も標数零では滑らかで、ここでの次元は不変接空間の18である。
これは完備正則局所環の余接空間に同変な持上げを選び、非自明固有文字のパラメータを消せば分かる。
微分階数の剰余検算を使う際、二十恒等式は局所化した整基底上で成り立つので全階数は180を超えず、
文字別の核も特殊化と可換である。
三例とも不変接ベクトルのうち、薄い各行列の零行・零列を一本も恒等的に零に保たないものがある。
このベクトルを固定軌跡で積分すると、一般点で薄い全行・全列が非零となる。
各薄い行列は階数1なので全成分が非零である。
全ての厚い交差ブロックの行列式も、基点で非零またはこの方向の一次項が非零である。
従って一般点では両交差ブロックが可逆で、階数2により対角ブロックも可逆である。
第二切断の色0の右上ブロックでは二つの対角積の和も非零で、薄い規格化の $z=-1$ を排除する。
その行列式の非零が $z=1$ を排除する。これで定理\ref{thm:klein-all-pairings}の全被覆に入る。
この固定成分は対角軌道の13次元を除いても5次元あり、各々4次元以下の他の四表示には入らない。
よって三点はいずれも五次元楕円族の閉包、従って既知の76次元成分に属する。
元の水平族は既約でそれぞれの点で全 $Z$ の成分が一意なので、族の全閉包について従う。
\end{proof}

\begin{proof}[垂直族での二つの別成分]
両型とも $b=c=1$, $t=\pm1$, $A=3t,B=1,D=7t$ とし、$C$ は76型で1、112型で4とする。
不変部分も含む全接空間の文字別次元の上界は正符号で $(17,18,20,20)$、負符号で $(17,16,16,16)$ である。
両側の基底変換の次元は先と同じで、射影安定群は有限である。

正符号では第四位置に文字1を加えると、$\mathrm{Sym}^2 E\otimes\mathrm{Sym}^2 F$ の元である。
この百次元空間の第一切断は $E\otimes F$ 上の対称行列であり、他方の隣接階数も同じになる。
対称な十六行行列に対する階数6以下の余次元は55なので、この固定した二同定の軌跡は各成分で少なくとも45次元である。
二同定を動かす自由度30を加えられる。実際に同じテンソルの同定の二通りから得られる二交換の合成は射影安定群に属し、
その有限性により同定を忘れるファイバーはここで有限である。
次元下界75と接次元上界75は等しく、全 $Z$ は滑らかで、通過する唯一の成分は一般に両対称となる。

同じ論証を四元群を固定して行う。第四位置の変更後の総文字は1であるため、両対称テンソルの空間は
$4\cdot2+2\cdot4+2\cdot2+2\cdot2=24$ 次元である。
二つずつの互いに転置な四色ブロックには階数2と1を一度ずつ課せばよく、余次元の上界は $4+9=13$。
各反対二位置の同定を同変な対角群内で動かす六自由度を加えると、不変軌跡も少なくとも17次元になる。
上の微分上界と一致して滑らかである。従ってこの軌跡の一般点でも、偶二位置の同定は文字を置換しない対角型、
奇二位置の同定は文字1の移動を伴う型である。
水平族と同じ薄い行列の微分検査により、一般点は全薄い成分が非零のチャートに達し、
$z=\pm1$ も除かれる。偶同定が対角型なので規格化では $M_2=N_2$, $M_3=N_3$ とできて $\lambda=1$。
全接次元が75の基点は76次元成分には含まれず、その楕円族側の閉包に入る可能性も消える。
残るのは $\lambda=1,A_1=A_0$ の既約な四変数表示である。
不変成分は対角軌道を除いて4次元なので、両型の正符号の基点はこの同じ表示の稠密開集合からの特殊化となる。
この表示の一般点で全 $Z$ の成分も一意である。実際、その閉包に含まれる上の滑らかな点の近傍は一般点を含む。
両型の正符号の垂直族の全閉包が一つの75次元成分に入る。

負符号は四元群二コサイクル捻りである。固定軌跡上でこの操作は線型な自己同型で、
正の垂直族とそれを含む四変数表示を、負の垂直族と $\lambda=-1$ の表示に移す。
不変軌跡の下界17から、全基底変換を許せば軌道の非自明文字部分の48次元をさらに加えられる。
厳密には正側の17次元不変軌跡が滑らかな点の像でこの48方向は不変接空間と交わらず、
あるいは有限射影安定群内の四元群の選び方が有限であることから次元下界65が従う。
負側の一般垂直点での上界65と一致し、四変数表示の基底変換による閉包そのものが全 $Z$ の既約成分になる。
全ての次元下界は同じ記号の標数零の点で得られ、剰余での接上界の等号もこれによって保証される。

一般同型類の各次元は一般軌道の61次元を引けばよい。垂直成分の両符号とも既知の正則点から
内部次数ごとの完全性を一般点へ持ち上げられる。水平・垂直の交叉点が特異であることは、
各々でこれら二つの異なる成分を通ることからも分かる。
\path{computations/verify_rectangle_component_smoothing.py} が点環の余法線、七つの代表の微分と薄い成分の変形、
および各型・符号の交叉例の文字別接次元 $(19,18,20,20)$ を検算する。
\end{proof}
